\begin{table}[htbp]
\caption{Sample Construction and Outcome Availability by Grade: Liberia}
\label{tab:lib_sampleflow}
\begin{center}
\begin{small}
\begin{threeparttable}
\begin{tabular}{l*{5}{r}}
\toprule
 & Grade 1 & Grade 2 & Grade 3 & Grade 4 & Total \\
\midrule
All observations & 1986 & 1937 & 2035 & 1627 & 7585 \\
Non-missing treatment & 1986 & 1937 & 2035 & 1627 & 7585 \\
Has baseline & 1184 & 1233 & 1243 & 1124 & 4784 \\
Final analytic sample & 1184 & 1233 & 1243 & 1124 & 4784 \\
\quad with endline & 839 & 840 & 745 & 730 & 3154 \\
\bottomrule
\end{tabular}
\begin{tablenotes}
\footnotesize \item Notes: The table reports student counts at each stage of sample construction for Liberia. Columns correspond to enrolled grade and the total column sums across grades. `All observations' is the raw student roster in the source files. `Has baseline' is the number of students with a non-missing baseline reading score used for assignment analysis. `Final analytic sample' is the baseline sample used in the main student-level specifications, and the indented row reports how many of those students also have a matched endline reading outcome.
\end{tablenotes}
\end{threeparttable}
\end{small}
\end{center}
\end{table}